\section{Introduction}\label{sec:introduction}
\IEEEPARstart{L}{ow-power} industrial Internet of Things (IIoT) deployments, such as process-monitoring and metering meshes, depend on the IPv6 Routing Protocol for Low-Power and Lossy Networks (RPL) \cite{rfc6550}. RPL builds a destination-oriented directed acyclic graph (DODAG) rooted at a border router using three control messages: DODAG Information Objects (DIOs) advertise the graph, DODAG Information Solicitations (DISs) request DIOs, and Destination Advertisement Objects (DAOs) build the downward routes that the root uses to reach every device. Because these messages are sent by the nodes themselves, an insider that holds valid network credentials can abuse them without breaking any cryptography \cite{albinali2024,raoof2019}.

Two such abuses are well documented. In multicast DIS flooding, a node solicits DIOs repeatedly, which resets the Trickle timers \cite{rfc6206} of its neighbors and floods the network with control traffic \cite{rajasekar2022,vermaranga2020,medjek2021}. In DAO manipulation, a node advertises a false transit parent, so the root installs a source route \cite{rfc6554} that does not exist and downward packets are lost \cite{pu2018,ghaleb2019,limsd2024}. Each has its own lightweight defense. DIS-FDM, for example, fuses DIS and parent-change evidence with Dempster--Shafer theory \cite{prajisha2025}, while Li-MSD counts DAOs to blacklist an insider \cite{limsd2024}. The problem this study addresses arises when one compromised node performs both attacks at the same time. A DIS-focused defense then sees the flood and may even isolate the attacker, but it never examines the DAO plane, so the false routes stay in the root's table and downward service remains broken. Detection and service separate. The authors of DIS-FDM themselves name this hybrid DIS and DAO scenario as future work \cite{prajisha2025}.

We call this threat a \emph{coordinated} attack in a narrow, same-node sense: each malicious node executes both mechanisms, and attackers do not need to exchange messages. Section~\ref{sec:related} shows that no reviewed defense has been evaluated against it, that DIS-FDM's belief values are fixed by hand, that its published evaluation reports no detection accuracy or statistical testing, and that no reviewed lightweight RPL defense lets each node learn its own response. These four gaps (Section~\ref{sec:gaps}) define the research.

We propose Adaptive Evidence Fusion for RPL (AEF-RPL), a single algorithm that runs on every node. It keeps DIS-FDM's evidence sources and adds four elements, each chosen for a stated reason. First, a DAO-consistency check compares the transit parent advertised in a DAO with the node that actually received it from its originator. In RPL, a DAO is sent upward to the preferred parent, so at the first hop the two must agree. Second, the uncertainty of this evidence is calibrated online from benign observations with a Beta-reputation estimate \cite{josang2002,ganeriwal2008} instead of being set by hand. Third, once a DAO sender is isolated, its DAOs are repaired rather than dropped, so that the routes of its honest descendants survive. Fourth, a tabular Q-learning layer \cite{watkins1992} decides whether to accept or isolate each neighbor from local reward, using 768 bytes of fixed-point state. A learning-disabled variant shares the same evidence, which lets us measure what learning adds on its own.

This study answers five research questions:
\begin{enumerate}
\item[\emph{RQ1:}] To what extent does an independent Contiki-NG/Cooja implementation reproduce the performance trends reported for DIS-FDM under multicast DIS flooding?
\item[\emph{RQ2:}] How much more does executing DIS flooding and DAO manipulation from the same nodes degrade delivery, control overhead, and resilience than either attack alone?
\item[\emph{RQ3:}] To what extent does adding self-calibrated DAO-consistency evidence and selective repair to DIS-FDM change detection accuracy and resilience against the combined attack?
\item[\emph{RQ4:}] How does the on-node Q-learning layer change false positives, detection latency, and resilience compared with the same detector using a fixed decision rule?
\item[\emph{RQ5:}] How consistent are these effects across seeds, attacker fractions, static and mobile topologies, attack rates, and radio reception loss?
\end{enumerate}
The questions follow the research proposal. Its RQ5 also covered networks of 500 and 1000 nodes; because the benign network at those sizes did not form stable routes in the simulator, attack rate and reception loss replaced network size as the further robustness conditions (Section~\ref{sec:limitations}).

The main contributions are as follows:
\begin{enumerate}
\item \emph{AEF-RPL}, a fully local defense against same-node DIS flooding and DAO manipulation. It combines a first-hop DAO-consistency invariant, Beta-calibrated evidence uncertainty, a two-window persistence rule, selective DAO repair, and an on-node Q-learning decision layer that fits in 768 bytes.
\item A \emph{mathematical analysis} that bounds the DAO-channel detection delay and false-alarm probability, shows how conflicting evidence slows the alarm, derives the learner's decision threshold, and proves that the fixed-point Q-values cannot overflow.
\item Two \emph{service-oriented metrics}, the Attack Resilience Index (ARI) and the Downward Recovery Time (DRT). They measure retained round-trip service against a matched attack-free run and so count nodes that lose their route, which conditional packet delivery ratios ignore.
\item A \emph{paired evaluation} of 2720 simulation runs against four published defenses, with pre-specified hypotheses, component ablations, and rate and radio-loss robustness studies. The evaluation shows where the gain comes from (selective DAO repair), what learning adds (faster decisions), and where the design falls short (false isolation under mobility).
\item A \emph{reproducibility artifact} with firmware, scenarios, and all 2720 raw logs, from which one command regenerates every table, test, and figure.
\end{enumerate}

The paper is organized as follows. Section~\ref{sec:related} reviews related work and states the research gaps. Section~\ref{sec:design} links the gaps to the objectives, hypotheses, and conceptual framework. Section~\ref{sec:method} presents AEF-RPL, and Section~\ref{sec:theory} analyzes it. Section~\ref{sec:setup} describes the experimental methodology, Section~\ref{sec:results} reports the results by hypothesis, and Section~\ref{sec:discussion} discusses them. Section~\ref{sec:conclusion} concludes.
